\documentclass[10pt,a4paper]{amsart}
\usepackage[margin=19mm]{geometry}
\usepackage{amsmath,amssymb,mathtools}
\usepackage[scr=rsfs,cal=euler]{mathalfa}
\usepackage{xfrac}
\usepackage[unicode,hidelinks,bookmarks=false]{hyperref}
\newcommand{\ie}{i.e.\ }
\newcommand{\cf}{cf.\ }
\newcommand{\resp}{respectively\ }
\newcommand{\Subsection}{\subsection}
\providecommand{\nobreakdash}{\nobreak\hbox{-}}
\setlength{\emergencystretch}{2em}
\multlinegap0pt
\usepackage{mdframed}
\usepackage{mdframed}
\usepackage{mdframed}
\usepackage{mdframed}
\usepackage{amssymb}
\usepackage{dsfont}
\usepackage{csquotes}
\usepackage{mdframed}
\usepackage[shortlabels]{enumitem}
\newtheorem{notation}{Notation}
\newtheorem{definition}{Definition}
\newtheorem{proposition}{Proposition}
\def\Sd{\mathbb{S}^d}
\def\T{\mathcal{Q}}
\def\X{\mathcal X}
\def\Zd{\mathbb{Z}^d} 
\def\beq{\begin{equation}}
\def\bz{\mathfrak{z}\hspace*{-4pt}\mathfrak{z}}
\def\eeq{\end{equation}}
\def\p{\mathfrak{p}}
\def\s{\mathfrak{s}}
\title{The fibre operators in the Bloch-Floquet decomposition\\of periodic magnetic pseudo-differential operators}
\author{Horia D. Cornean, Bernard Helffer, Radu Purice}
\date{Source: arXiv:2512.22547v2 (CC BY 4.0)}
\begin{document}
\maketitle

\noindent\emph{Source and licence.} The fibre operators in the Bloch-Floquet decomposition\\of periodic magnetic pseudo-differential operators. Version arXiv:2512.22547v2 (CC BY 4.0), distributed by arXiv under the Creative Commons Attribution 4.0 International licence, http://creativecommons.org/licenses/by/4.0/, as recorded in the arXiv licence metadata for this exact version and shown on the abstract page https://arxiv.org/abs/2512.22547v2. Full licence notice: \texttt{licenses/arxiv-2512.22547-CC-BY-4.0.txt}.\par\smallskip

\noindent\textit{Editorial scope.} Counted unit: the article's proposition F-comp-torWsyst (source0.tex lines 1136--1144). The article's complete original proof of it is delivered with it (source0.tex lines 1145--1165, one \texttt{proof} environment of 21 lines). No mathematics is written, repaired or re-derived: the statement and the proof are the article's own lines, in the article's own notation, and no retained line is substituted or re-flowed.  Retained uncounted as the source-local context the proof needs: lines 1121--1128; lines 1130--1135.  Nothing was omitted from any retained span, so the package carries no omission record.  No retained line contains a citation, so the wrapper carries no bibliography and no bibliography entry.  No retained line contains a figure environment, an image-inclusion command or drawing code, so no image is delivered and the package declares no asset dependency.  Editorial formatting only, in addition: an AMS \texttt{amsart} preamble that restates the article's own declarations and macros used by the retained lines, namely the environments \texttt{notation}, \texttt{definition}, \texttt{proposition} on separate counters, together with the standard packages those lines need.  The retained environments print in this document as Notation 1, Definition 1, Proposition 1 in order of appearance, whereas the article prints the same items with the numbering of its own full document.  The complete native source of the article as distributed in the arXiv e-print for this exact version is bundled unchanged as \texttt{sources/191743/source0.tex}.
\begin{notation}
Let us denote by $\T_2:=\X/[2\Gamma]$ that is  isomorphic with $\Sd_2$.
\end{notation}
As the two Pontryagin duals of $\T$ and $\T_2$ are isomorphic with $\Zd$, in the following arguments we shall use the following two isomorphisms:
\beq\label{DF-char}\begin{split}
	&\Zd\ni\nu\mapsto\hat{\chi}_{1,\nu}\in\widehat{\Sd},\quad\hat{\chi}_{1,\nu}(\bz):=\bz^{-\nu}=e^{-i<\nu,\s(\bz)>}\,,\\
	&\Zd\ni\nu\mapsto\hat{\chi}_{2,\nu}\in\widehat{\Sd_2},\quad\hat{\chi}_{2,\nu}(\tilde{\bz}):=\tilde{\bz}^{-\nu}=e^{-i<(1/2)\nu,\s_2(\tilde{\bz})>}\,.
\end{split}\eeq

\begin{definition}\label{D-s-tor-W-syst}
We call   \emph{symmetric Weyl representation of  $\T\times\Gamma_*$} the unitary valued smooth map (see also \eqref{DF-char} for the notations):
	\beq\label{DF-s-tor-W-syst}
	 \T_2\times\Gamma_*\ni\,(\tilde{\bz},\gamma^*)\,\mapsto\,\widetilde{\mathbb{W}}_{\T}(\tilde{\bz},\gamma^*):=\hat{\chi}_{2,-\gamma^*}(\tilde{\bz})U_{\T}\big(\tilde{\p}_2(\tilde{\bz})\big)V_{\Gamma_*}(\gamma^*)\in\mathbb{U}\big(L^2(\T)\big).
	\eeq
\end{definition}

\begin{proposition}
The symmetric Weyl system { $\widetilde{\mathbb{W}}_{\T}:\T_2\times\Gamma_*\rightarrow\mathbb{U}\big(L^2(\T)\big)$ given in Definition \ref{D-s-tor-W-syst}} has the following properties:
\beq
\widetilde{\mathbb{W}}_{\T}(\tilde{\bz},\gamma^*)^*=\widetilde{\mathbb{W}}_{\T}(\tilde{\bz}^{-1},-\gamma^*);
\eeq
\beq\label{F-comp-torWsyst}
\widetilde{\mathbb{W}}_{\T}(\tilde{\bz},\alpha^*)\widetilde{\mathbb{W}}_{\T}(\tilde{\bz'},\beta^*)=e^{(i/2)[<\alpha^*,\s_2(\tilde{\bz}')>-<\beta^*,\s_2(\tilde{\bz})>]}\widetilde{\mathbb{W}}_{\T}(\tilde{\bz}\tilde{\bz'},\alpha^*+\beta^*).
\eeq
\end{proposition}

\begin{proof}
We have the following equalities:
\begin{align*}
	\widetilde{\mathbb{W}}_{\T}(\tilde{\bz},\gamma^*)^*&=\hat{\chi}_{2,\gamma^*}(\tilde{\bz})V_{\Gamma_*}(\gamma^*)^*U_{\T}(\tilde{\bz}^{2})^*=\hat{\chi}_{2,\gamma^*}(\tilde{\bz})\hat{\chi}_{1,-\gamma^*}(\tilde{\bz}^{2})^*U_{\T}(\tilde{\bz}^{2})^*V_{\Gamma_*}(\gamma^*)^*\\
	&=\exp\big((i/2)<\gamma^*,\s_2(\tilde{\bz})>-i<\gamma^*,\s(\tilde{\bz}^2)>\big)U_{\T}(\tilde{\bz}^{-2})V_{\Gamma^*}(-\gamma^*),
\end{align*}
\begin{align*}
e^{(i/2)<\gamma^*,\s_2(\tilde{\bz})>-i<\gamma^*,\s(\tilde{\bz}^2)>}&=e^{(i/2)<\gamma^*,\s(\tilde{\bz}^2)+\kappa(\tilde{\bz})>-i<\gamma^*,\s(\tilde{\bz}^2)>}\\
&=e^{-(i/2)<\gamma^*,\s(\tilde{\bz}^2)>+(i/2)<\gamma^*,\kappa(\tilde{\bz})>}\\
&=e^{-(i/2)<\gamma^*,\s(\tilde{\bz}^2)>-(i/2)<\gamma^*,\kappa(\tilde{\bz})>}\\
&=e^{-(i/2)<\gamma^*,\s_2(\tilde{\bz})>}=\hat{\chi}_{2,\gamma^*}(\tilde{\bz}^{-1}).
\end{align*}
Thus, we have $$\widetilde{\mathbb{W}}_{\T}(\tilde{\bz},\gamma^*)^*=\widetilde{\mathbb{W}}_{\T}(\tilde{\bz}^{-1},-\gamma^*)\,.$$

We also obtain that:
\begin{align*}
	\widetilde{\mathbb{W}}_{\T}(\tilde{\bz},\alpha^*)\tilde{\mathbb{W}}_{\T}(\tilde{\bz'},\beta^*)&=\hat{\chi}_{2,-\alpha^*}(\tilde{\bz})U_{\T}(\tilde{\bz}^2)V_{\Gamma_*}(\alpha^*)\hat{\chi}_{2,-\beta^*}(\tilde{\bz'})U_{\T}(\tilde{\bz'}^2)V_{\Gamma_*}(\beta^*)\\ \nonumber
&=e^{i<\alpha^*,\s(\tilde{\bz'}^2)>}\,e^{(i/2)<\alpha^*,\s_2(\tilde{\bz})>}\,e^{(i/2)<\beta^*,\s_2(\tilde{\bz}')>}\,e^{-(i/2)<\alpha^*+\beta^*,\s_2(\tilde{\bz}\,\tilde{\bz'})>}\widetilde{\mathbb{W}}_{\T}(\tilde{\bz}\tilde{\bz'},\alpha^*+\beta^*)\\ \nonumber
	&=e^{(i/2)[<\alpha^*,\s_2(\tilde{\bz}')>-<\beta^*,\s_2(\tilde{\bz})>]}\widetilde{\mathbb{W}}_{\T}(\tilde{\bz}\tilde{\bz'},\alpha^*+\beta^*).
\end{align*}
\end{proof}

\end{document}
